% sec_strip_lattice.tex -- integer-orbit projection for the observation channel
\section{Integer-Orbit Projection for the Strip RKHS}%
\label{sec:strip-lattice}

This appendix records only the strip facts needed by the observation channel
from Theorem~\ref{thm:poisson-channel-rkhs}.  The kernel is
\[
K(a,b):=\frac{2}{4+(a-b)^2}=\pi P_2(a-b),
\qquad a,b\in\RR,
\]
and $\mathcal H_K$ denotes its complex reproducing-kernel Hilbert space.  With
the Fourier convention of this paper,
\[
\langle f,g\rangle_{\mathcal H_K}
=\frac{1}{2\pi^2}\int_{\RR}e^{2|\xi|}\widehat f(\xi)
\overline{\widehat g(\xi)}\,\dd\xi,
\]
so $\mathcal H_K$ is the Poisson image space
\[
\mathcal H_K
=\left\{f:\int_{\RR}e^{2|\xi|}|\widehat f(\xi)|^2\dd\xi<\infty\right\}
\]
with this norm.  This is the standard translation-invariant RKHS
realization of a positive definite kernel, since
\[
\widehat{K(\cdot,0)}(\xi)=\pi e^{-2|\xi|}.
\]
Every $f\in\mathcal H_K$ extends holomorphically to the strip
$\{z\in\CC:|\Im z|<1\}$, but the lattice argument below uses only the real-line
Fourier model and the reproducing property.

The entropy interface is isolated in
Theorem~\ref{thm:observation-entropy-interface}.  The entropy-defect results
use only a finite centred moment jet.  The strip observation channel supplies
that jet when the full profile is available, either directly or from integer
samples in the exact lattice case.  Integer samples without the lattice
condition are projection data only and are not used below as coefficient or
defect data.  The projection theorem in this appendix is not used to prove the
entropy coefficient identities; it is a companion observation-channel
interface for translating separately supplied full-profile data into the
moment inputs that those identities require.

\begin{lemma}[Integer projection for the Poisson strip kernel]%
\label{lem:auxiliary-uniform-lattice-projection}
\label{prop:poisson-lattice-symbol}
Let
\[
\mathcal H_K(\ZZ)
:=
\overline{\operatorname{span}}\{K(\cdot,n):n\in\ZZ\}
\subset\mathcal H_K
\]
and set
\[
\Sigma_1(\theta):=\sum_{m\in\ZZ}e^{-2|\theta+2\pi m|},
\qquad -\pi\le\theta\le\pi .
\]
Then
\begin{equation}\label{eq:lattice-symbol}
\sigma(\theta):=\sum_{n\in\ZZ}K(n,0)e^{-\ii n\theta}
=\pi\Sigma_1(\theta)
=\pi\frac{\cosh(2(\pi-|\theta|))}{\sinh(2\pi)}
\qquad(-\pi\le\theta\le\pi).
\end{equation}
Moreover, if
\[
q_f(\xi):=e^{2|\xi|}\widehat f(\xi),
\]
then $f\in\mathcal H_K(\ZZ)$ if and only if $q_f$ has a $2\pi$-periodic
representative satisfying
\[
\int_{-\pi}^{\pi}\Sigma_1(\theta)|q_f(\theta)|^2\,\dd\theta<\infty.
\]
In that case
\begin{equation}\label{eq:lattice-subspace-norm}
\|f\|_{\mathcal H_K}^2
=\frac{1}{2\pi^2}\int_{-\pi}^{\pi}
\Sigma_1(\theta)|q_f(\theta)|^2\,\dd\theta .
\end{equation}
Let $\Pi_{\ZZ}$ be the orthogonal projection of $\mathcal H_K$ onto
$\mathcal H_K(\ZZ)$.  For every $f\in\mathcal H_K$,
\begin{equation}\label{eq:integer-orbit-projection-symbol}
q_{\Pi_{\ZZ}f}(\theta)
=
\frac{\displaystyle\sum_{m\in\ZZ}\widehat f(\theta+2\pi m)}
{\displaystyle\sum_{m\in\ZZ}e^{-2|\theta+2\pi m|}},
\qquad -\pi\le\theta\le\pi,
\end{equation}
extended $2\pi$-periodically.  The alias sum is absolutely convergent for
almost every $\theta$, and
\[
\widehat{\Pi_{\ZZ}f}(\xi)
=e^{-2|\xi|}q_{\Pi_{\ZZ}f}(\xi).
\]
The integer samples determine precisely this projection:
\[
f(n)=g(n)\quad(n\in\ZZ)
\qquad\Longleftrightarrow\qquad
\Pi_{\ZZ}f=\Pi_{\ZZ}g .
\]
\end{lemma}

\begin{proof}
Poisson summation applied to $k(x)=K(x,0)$ gives
\[
\sigma(\theta)
=\sum_{m\in\ZZ}\widehat k(\theta+2\pi m)
=\pi\sum_{m\in\ZZ}e^{-2|\theta+2\pi m|}.
\]
For $|\theta|\le\pi$ the last series is
\[
e^{-2|\theta|}
+\sum_{m=1}^\infty e^{-2(|\theta|+2\pi m)}
+\sum_{m=1}^\infty e^{-2(2\pi m-|\theta|)}
=\frac{\cosh(2(\pi-|\theta|))}{\sinh(2\pi)},
\]
which proves \eqref{eq:lattice-symbol}.

If $F=\sum_{|n|\le N}c_nK(\cdot,n)$, then
\[
\widehat F(\xi)=\pi e^{-2|\xi|}\sum_{|n|\le N}c_ne^{-\ii n\xi},
\qquad
q_F(\xi)=\pi\sum_{|n|\le N}c_ne^{-\ii n\xi}.
\]
Thus $q_F$ is $2\pi$-periodic, and the norm formula follows by partitioning
$\RR$ into the translates $[-\pi,\pi]+2\pi m$:
\[
\|F\|_{\mathcal H_K}^2
=\frac{1}{2\pi^2}\int_{-\pi}^{\pi}
\Sigma_1(\theta)|q_F(\theta)|^2\,\dd\theta.
\]
Taking closures proves one implication.  Conversely, if $q$ is periodic with
finite weighted norm, trigonometric polynomials are dense in the equivalent
space $L^2([-\pi,\pi],\Sigma_1(\theta)\dd\theta)$, because $\Sigma_1$ is
continuous and strictly positive.  Approximating $q/\pi$ by such polynomials
and using the displayed norm identity gives an element of
$\mathcal H_K(\ZZ)$ with profile $q$.  This proves the characterization of
$\mathcal H_K(\ZZ)$ and \eqref{eq:lattice-subspace-norm}.

For $f\in\mathcal H_K$, Cauchy--Schwarz gives, for almost every
$\theta\in[-\pi,\pi]$,
\[
\sum_{m\in\ZZ}|\widehat f(\theta+2\pi m)|
\le
\Sigma_1(\theta)^{1/2}
\left(\sum_{m\in\ZZ}e^{2|\theta+2\pi m|}
|\widehat f(\theta+2\pi m)|^2\right)^{1/2}.
\]
Hence the alias sum in \eqref{eq:integer-orbit-projection-symbol} is defined
almost everywhere.  For $g\in\mathcal H_K(\ZZ)$, with periodic profile $q_g$,
the norm formula on the fibers gives
\[
\|f-g\|_{\mathcal H_K}^2
=\frac{1}{2\pi^2}\int_{-\pi}^{\pi}
\sum_{m\in\ZZ}e^{-2|\theta+2\pi m|}
|q_f(\theta+2\pi m)-q_g(\theta)|^2\,\dd\theta.
\]
For each fiber, this quadratic expression is minimized by the weighted
average
\[
q_g(\theta)
=\frac{\sum_m e^{-2|\theta+2\pi m|}q_f(\theta+2\pi m)}
{\Sigma_1(\theta)}
=\frac{\sum_m\widehat f(\theta+2\pi m)}{\Sigma_1(\theta)}.
\]
This is the profile of the orthogonal projection and proves
\eqref{eq:integer-orbit-projection-symbol}.

Finally, the reproducing property gives $f(n)=\langle f,K(\cdot,n)\rangle$.
Since $\mathcal H_K(\ZZ)$ is the closed span of these kernel sections,
$f-g$ vanishes on all integers exactly when $f-g\perp\mathcal H_K(\ZZ)$.
This is equivalent to $\Pi_{\ZZ}f=\Pi_{\ZZ}g$.
\end{proof}

\begin{observationprofilelemma}
\label{lem:observation-profile-law-inversion}
Let $\nu$ be a probability measure on $\RR$ with finite mean $\bar\gamma$,
let \(t>0\), and set
\[
\Delta:=\gamma-\bar\gamma,\qquad
\mu_t:=\operatorname{Law}(\Delta/t).
\]
Define the finite signed measure
\[
\lambda_t(\dd s):=s\,\mu_t(\dd s).
\]
For the Poisson observation channel \(F_t\) of
Theorem~\ref{thm:poisson-observation-lattice-sampling}, with Fourier
convention \(\widehat\lambda_t(\xi)=\int_{\RR}e^{-\ii s\xi}\dd\lambda_t(s)\),
\begin{equation}\label{eq:observation-profile-fourier-inversion}
q_{F_t}(\xi)=\pi\widehat{\lambda_t}(\xi).
\end{equation}
Consequently \(q_{F_t}\) determines \(\lambda_t\) by Fourier uniqueness for
finite signed measures.  It determines \(\mu_t\) on \(\RR\setminus\{0\}\) by
\begin{equation}\label{eq:observation-profile-offzero-recovery}
\mu_t(B)=\int_B s^{-1}\dd\lambda_t(s),
\qquad 0\notin\overline B,
\end{equation}
and then determines the remaining atom by
\begin{equation}\label{eq:observation-profile-zero-atom-recovery}
\mu_t(\{0\})=1-\mu_t(\RR\setminus\{0\}).
\end{equation}
Thus full knowledge of \(q_{F_t}\) determines the signed measure
\(s\,\mu_t(\dd s)\) and then the full scaled centred law \(\mu_t\); this is a
measure-inversion statement, separate from any integer-sampling statement.
\end{observationprofilelemma}

\begin{proof}
The identity follows from
\[
\frac{\pi}{t}\int_{\RR}\Delta e^{-\ii\Delta\xi/t}\dd\nu(\gamma)
=\pi\int_{\RR}s e^{-\ii s\xi}\dd\mu_t(s).
\]
The right-hand side is \(\pi\widehat{\lambda_t}(\xi)\).  Since
\(\lambda_t\) has finite total variation, its Fourier transform determines
\(\lambda_t\) uniquely.  Formula
\eqref{eq:observation-profile-offzero-recovery} is division of the signed
measure \(\lambda_t=s\mu_t\) by the bounded continuous function \(s\) on sets
bounded away from zero.  A monotone exhaustion of \(\RR\setminus\{0\}\)
recovers the whole off-zero restriction of \(\mu_t\), and
\eqref{eq:observation-profile-zero-atom-recovery} uses
\(\mu_t(\RR)=1\).
\end{proof}

\begin{theorem}[Poisson observation-channel interface]%
\label{thm:poisson-observation-lattice-sampling}
\label{prop:poisson-law-lattice-periodicity}
Let $\nu$ be a probability measure on $\RR$ with finite mean $\bar\gamma$,
let $t>0$, and put
\[
\Delta:=\gamma-\bar\gamma,
\qquad
\mu_t:=\operatorname{Law}(\Delta/t).
\]
Let
\[
h_t:=P_t*\nu,\qquad g_t(x):=P_t(x-\bar\gamma).
\]
Define
\[
d_t(y):=
\frac{h_t(\bar\gamma+ty)}{g_t(\bar\gamma+ty)}-1 .
\]
Then \(d_t\in L^2_0(\omega)\).  Writing \(U:=\mathcal U\), the unitary
\(U:L^2_0(\omega)\to\mathcal H_K\) gives
\[
F_t:=Ud_t\in\mathcal H_K,
\]
and
\begin{equation}\label{eq:poisson-law-side-q-profile}
q_{F_t}(\xi):=e^{2|\xi|}\widehat F_t(\xi)
=\pi\int_{\RR}s e^{-\ii s\xi}\,\mu_t(\dd s)
\qquad\text{for a.e. }\xi .
\end{equation}
Equivalently, if \(\mathfrak d_t\) denotes the centred fluctuation profile of
Theorem~\ref{thm:poisson-channel-rkhs}, then \(d_t=\mathfrak d_t\) and
the integer observation samples are
\begin{equation}\label{eq:poisson-observation-integer-samples}
F_t(n)
=\langle d_t,\Psi_n\rangle_{L^2(\omega)}
=\frac1t\int_{\RR}\Delta\,
K\!\left(n,\frac{\Delta}{t}\right)\dd\nu(\gamma),
\qquad n\in\ZZ .
\end{equation}
The recovery conclusions are as follows.
\begin{enumerate}[label=\textup{(\roman*)},leftmargin=*]
\item The integer samples recover \(\Pi_{\ZZ}F_t\) and only
\(\Pi_{\ZZ}F_t\) among ambient \(\mathcal H_K\) data.
\item They recover \(F_t\) exactly if and only if
\(\mu_t\) is supported on \(\ZZ\), equivalently
\begin{equation}\label{eq:poisson-law-side-lattice-support}
\nu\bigl(\{\gamma\in\RR:\gamma-\bar\gamma\in t\ZZ\}\bigr)=1 .
\end{equation}
\item Once the full profile \(q_{F_t}\) is known, the separate
Observation-profile inversion lemma recovers the full scaled centred law
\(\mu_t\).
\end{enumerate}
The projection is obtained by inserting $\widehat F_t$ into
\eqref{eq:integer-orbit-projection-symbol}.

The integer samples recover the full channel $F_t$ exactly if and only if
\[
F_t\in\mathcal H_K(\ZZ)
\]
if and only if \(q_{F_t}\) has a $2\pi$-periodic representative.
In this exact-recovery case,
\begin{equation}\label{eq:poisson-law-side-lattice-norm}
\|F_t\|_{\mathcal H_K}^2
=\frac{1}{2\pi^2}\int_{-\pi}^{\pi}
\Sigma_1(\theta)|q_{F_t}(\theta)|^2\,\dd\theta .
\end{equation}
Finally, if the full profile \(q_{F_t}\) is available, whether from the
lattice-periodic integer data above or by any other observation of the full
channel, Lemma~\ref{lem:observation-profile-law-inversion} recovers
\(s\,\mu_t(\dd s)\), then \(\mu_t\) on \(\RR\setminus\{0\}\), and finally
\(\mu_t(\{0\})\) from total mass.
\end{theorem}

\begin{proof}
Theorem~\ref{thm:poisson-channel-rkhs} gives
\[
F_t(a)=\frac1t\int_{\RR}\Delta
K\!\left(a,\frac{\Delta}{t}\right)\dd\nu(\gamma).
\]
It also gives \(d_t=\mathfrak d_t\in L^2_0(\omega)\) and
\(F_t=\mathcal Ud_t\).  Taking $a=n$ and using
\(\mathcal U\Psi_n=K(\cdot,n)\) gives
\eqref{eq:poisson-observation-integer-samples}.  By
Lemma~\ref{lem:auxiliary-uniform-lattice-projection}, the integer samples of
an ambient function in $\mathcal H_K$ determine exactly its orthogonal
projection onto $\mathcal H_K(\ZZ)$; applying this to $F_t$ proves the
projection assertion.

The Fourier transform of the kernel in the first variable is
\[
\widehat{K(\cdot,b)}(\xi)=\pi e^{-2|\xi|}e^{-\ii b\xi}.
\]
Taking the Fourier transform under the Bochner integral for $F_t$ gives
\[
\widehat F_t(\xi)
=\frac{\pi}{t}e^{-2|\xi|}
\int_{\RR}\Delta e^{-\ii\Delta\xi/t}\dd\nu(\gamma),
\]
which is \eqref{eq:poisson-law-side-q-profile}.

By the lemma, $F_t\in\mathcal H_K(\ZZ)$ if and only if $q_{F_t}$ has a
$2\pi$-periodic representative.  The displayed representative is continuous
by dominated convergence and the finite first moment hypothesis, so a.e.
periodicity is the same as everywhere periodicity for this representative.

It remains to translate full-channel recovery into the law-side lattice
condition.  Let \(\mu=\mu_t\) be the law of \(\Delta/t\) and set
\(\lambda(\dd s):=s\,\mu(\dd s)\).
The periodicity of \(q_{F_t}\) is equivalent to
\[
\widehat\lambda(\xi+2\pi)=\widehat\lambda(\xi)
\qquad(\xi\in\RR).
\]
Thus the finite signed measure $(e^{-\ii2\pi s}-1)\lambda(\dd s)$ has
identically zero Fourier transform.  Uniqueness of Fourier transforms of
finite measures implies
\[
(e^{-\ii2\pi s}-1)s\,\mu(\dd s)=0.
\]
Hence $\mu$ is supported on $\ZZ$; the point $s=0$ is already in $\ZZ$.
This is exactly \eqref{eq:poisson-law-side-lattice-support}.  Conversely, if
$\Delta/t\in\ZZ$ almost surely, then $e^{-\ii2\pi\Delta/t}=1$ almost surely,
so \(q_{F_t}\) is $2\pi$-periodic.  The norm formula is then
\eqref{eq:lattice-subspace-norm} applied to $f=F_t$.

The final law-recovery assertion is precisely
Lemma~\ref{lem:observation-profile-law-inversion}: full channel knowledge gives
the full profile \(q_{F_t}\), which gives \(\lambda=s\mu_t\) by Fourier
uniqueness, gives \(\mu_t\) off zero by division by \(s\), and gives the atom at
zero from \(\mu_t(\RR)=1\).
\end{proof}

\begin{theorem}[Observation-to-entropy coefficient interface]%
\label{thm:observation-entropy-interface}
Let \(\nu\) be a probability measure on \(\RR\) with mean \(\bar\gamma\),
positive variance \(\sigma^2\), and finite centred sixth absolute moment.
Fix \(t_0>0\), put
\[
\Delta:=\gamma-\bar\gamma,\qquad
\mu_{t_0}:=\operatorname{Law}(\Delta/t_0),
\]
and let \(F_{t_0}\) be the Poisson observation channel from
Theorem~\ref{thm:poisson-observation-lattice-sampling}.

The following observation data determine the centred moment jet
\((\mu_2,\ldots,\mu_6)\), where
\(\mu_j=\int_{\RR}\Delta^j\,\dd\nu(\gamma)\):
\begin{enumerate}[label=\textup{(\roman*)},leftmargin=*]
\item the full continuous profile
\(q_{F_{t_0}}(\xi)=e^{2|\xi|}\widehat F_{t_0}(\xi)\) on \(\RR\);
\item the integer samples \((F_{t_0}(n))_{n\in\ZZ}\), but only together with
the exact-recovery condition
\[
\nu\{\gamma:\gamma-\bar\gamma\in t_0\ZZ\}=1 .
\]
\end{enumerate}
In either case the moments are recovered from the scaled law by
\begin{equation}\label{eq:observation-interface-moments}
\mu_j=t_0^j\int_{\RR}s^j\,\dd\mu_{t_0}(s),
\qquad 2\le j\le6.
\end{equation}
Consequently these observation data determine the coefficient quantities used
later in the entropy-defect results:
\[
C_6^{\rm ent}(\nu),\qquad A_8^{\rm ent}(\nu),\qquad
\Delta_6^{\rm ent}(\nu),\qquad \Delta_8^{\rm ent}(\nu),
\]
as defined and identified in Theorem~\ref{thm:poisson-kl-eighth} and in the
defect corollaries of Section~\ref{subsec:eighth-order-defects}.

If only the integer samples \((F_{t_0}(n))_{n\in\ZZ}\) are given without the
lattice-support condition, the data determine only the projection
\(\Pi_{\ZZ}F_{t_0}\).  No entropy-defect theorem in this paper uses
\(\Pi_{\ZZ}F_{t_0}\) alone to infer the centred moment jet, the coefficients
\(C_6^{\rm ent}\), \(A_8^{\rm ent}\), or the defects.  The finite-time
oracle estimate in
Theorem~\ref{thm:finite-time-poisson-entropy-stability} also requires its
separate eighth-order remainder inputs
\(B_{\nu,8}\) and \(T_{\nu,8}\), or their shell-uniform replacements; these
are not supplied by integer projection data alone.
\end{theorem}

\begin{proof}
Full profile data determine \(\mu_{t_0}\) by
Lemma~\ref{lem:observation-profile-law-inversion}: first recover the signed
measure \(s\,\mu_{t_0}(\dd s)\), then divide by \(s\) off the origin, and
finally recover the atom at the origin from total mass.  Since \(\nu\) has a
finite centred sixth absolute moment, the moments in
\eqref{eq:observation-interface-moments} are finite and give precisely the
centred moments of \(\nu\) through order six.

If integer samples are accompanied by the lattice-support condition, then
Theorem~\ref{thm:poisson-observation-lattice-sampling} says that the
integer-orbit reconstruction is the full channel \(F_{t_0}\).  This reduces
the second data set to the full-profile case.

The formulas for \(C_6^{\rm ent}\) and \(A_8^{\rm ent}\) in
Theorem~\ref{thm:poisson-kl-eighth} are polynomials in
\(\mu_2,\ldots,\mu_6\), and the defect variables in
Corollaries~\ref{thm:poisson-defect-ladder}--%
\ref{thm:poisson-defect-amplification} are defined from those identified
entropy coefficients and \(\sigma^2=\mu_2\).  Hence the recovered moment jet
determines exactly the listed entropy-defect quantities.

Without the lattice-support condition, the projection theorem gives only
\(\Pi_{\ZZ}F_{t_0}\).  The theorem therefore makes no coefficient-recovery
assertion from such projected data.  The statement about the finite-time
oracle inputs is the input separation made in
Theorem~\ref{thm:finite-time-poisson-entropy-stability}.
\end{proof}
